\section{Technical Advancements}\label{sec:technical_advancements}

\subsection{Small Language Models at Edges}\label{subsec:small_language_models}


The first move of AI to Edge is to deploy small language models (SLMs) to edge devices \cite{lu2024small,wang2024comprehensive,van2024survey}. This trend is driven by the growing demand for AI applications that can run directly on edge devices, motivated by needs for privacy, offline usage, and real-time processing without cloud dependence. However, edge devices have limited memory, computation, and energy resources, requiring more efficient and compact models.


\paragraph{SLMs leverage innovative architectures for efficient edge deployment.}
The classic Transformer architecture \cite{vaswani2017attention} uses self-attention mechanisms for effective sequence modeling but faces quadratic complexity challenges, with models like TinyBERT \cite{jiao2020tinybert} (14.5M parameters) and ALBERT \cite{lan2020albert} (12M parameters) demonstrating its effectiveness at small scales. Mamba \cite{gu2023mamba}, based on state space models, achieves linear complexity and faster inference by utilizing only the previous hidden state, as demonstrated by Zamba2-2.7B \cite{glorioso2024zamba2} which achieves twice the speed and 27\% reduced memory overhead compared to traditional models. Hymba \cite{dong2024hymba} combines both approaches by integrating attention and SSM heads within the same layer for parallel processing, with its 1.5B variant trained on DCLM-Baseline-1.0 and SmolLM-Corpus achieving 11.67 times cache size reduction while outperforming Llama-3.2-3B. The xLSTM architecture \cite{beck2024xlstm} modernizes LSTM with exponential gates and matrix memory cells, with models ranging from 125M to 1.3B parameters trained on 300 billion tokens from SlimPajama \cite{cerebras2023slimpajama}, consistently outperforming comparable RWKV-4 \cite{peng-etal-2023-rwkv}, Llama \cite{touvron2023llama3}, and Mamba models across various tasks in the PALOMA benchmark \cite{magnusson2024paloma}. These architectural innovations demonstrate the potential for efficient and powerful language models that can run effectively on edge devices.


\paragraph{SLMs can be constructed through diverse methodological approaches.}
The construction of efficient SLMs relies on a comprehensive suite of techniques, each with specific performance trade-offs. 
For training SLMs from scratch, optimized MLM approaches \cite{devlin2019bert} with increased masking ratio (25\% vs traditional 15\%) demonstrate 2-3\% performance improvements for models under 3B parameters.
When deriving SLMs from existing LLMs, knowledge distillation has proven particularly effective, with response-based distillation \cite{sanh2019distilbert} reducing model size by 40\% while maintaining 95\% of the original performance.
In architecture optimization, the Mixture of Experts approach \cite{fedus2022switch} enables 65\% parameter reduction while potentially improving performance in specific tasks. Domain specialization has shown remarkable results, particularly in the medical field where 3B parameter models achieve 92\% accuracy, outperforming 175B models (89\%) \cite{fu2023specializing}.
The combination of these techniques yields impressive results - a notable example is a 770M parameter model from \cite{hsieh2023distilling} that combines distillation, quantization, and domain specialization to achieve 95\% of the performance of a 540B model on specific tasks while requiring less than 0.15\% of the computational resources. Most successful SLMs achieve their efficiency by combining multiple techniques, typically starting with knowledge distillation, applying compression methods, and finishing with domain-specific fine-tuning. 
\cite{wang2024comprehensive} has provided a comprehensive survey of SLMs, and we summarize the architecture innovations and training methods in Appendix \ref{app:slm_architectures_training}. \looseness-1


Despite remarkable progress in deploying compressed models to edge devices, the current landscape remains largely confined to individual devices operating in isolation, failing to leverage the massive distributed computing power that could be achieved through collaborative training across edge devices. This represents a significant missed opportunity, as the collective computing resources of billions of edge devices worldwide remain untapped, while individual devices struggle with the computational demands of modern AI applications.

\subsection{Collaborative Inference at Edges}\label{subsec:collaborative_inference}
The emergence of collaborative inference at the edge represents a significant shift in AI infrastructure, enabling more accessible and cost-effective AI solutions compared to traditional data centers which often present barriers in terms of cost, energy consumption, and accessibility. Recent frameworks like Exo \cite{exo} enable users to create AI clusters using everyday devices such as phones, tablets, and computers through peer-to-peer architecture, effectively unifying their computational resources.

Several approaches advance this paradigm: Neurosurgeon \cite{kang2017neurosurgeon} introduces a lightweight scheduler that partitions DNN computations between devices and datacenters; MoE$^2$ \cite{jin2025moe2} optimizes LLM inference under energy and latency constraints; Edgent \cite{li2018edgent} enables low-latency edge intelligence through adaptive partitioning and early-exit mechanisms; and Galaxy \cite{ye2024galaxy} leverages hybrid model parallelism to efficiently execute Transformer models across edge devices. Despite these advances, the current landscape has yet to fully capitalize on the potential of distributed data resources at the edge. The next frontier lies in transforming these devices from mere inference endpoints into active training participants, representing a significant opportunity for distributed AI development. \looseness-1

\subsection{Feasibility of On-Device Training}

Recent advancements in optimization techniques have significantly reduced the memory and computational requirements for training machine learning models, making on-device training increasingly feasible, even on resource-constrained edge devices.
For instance, Lin et al. \cite{lin2022ondevice} demonstrated training neural networks on microcontroller units with only 256KB of RAM by employing an algorithm-system co-design framework. 
% Their approach includes Quantization-Aware Scaling to stabilize 8-bit quantized training and Sparse Update to skip gradient computation of less important layers, making it particularly suitable for resource-constrained environments.
Similarly, Cai et al. \cite{cai2020tinytl} introduced a tiny transfer learning approach that freezes most parameters and only trains a small subset, allowing effective fine-tuning with minimal memory requirements. Qiu et al. \cite{qiu2022zerofl} proposed ZeroFL, a framework that relies on highly sparse operations to accelerate on-device training in federated learning settings, enabling efficient model training on edge devices with up to 95\% sparsity.
Recent work by Sugiura and Matsutani \cite{sugiura2025elasticzo} further advanced this field with ElasticZO, which combines zeroth-order and first-order optimization to achieve a better trade-off between accuracy and training cost. Their ElasticZO-INT8 variant achieves integer arithmetic-only training, further reducing memory usage and training time by approximately 1.5x without compromising accuracy.
% Zeroth-order methods have also been extended to federated learning for large language models as demonstrated by Ling et al. \cite{ling2024convergence}, who showed that memory-efficient zeroth-order optimization can reduce GPU memory usage during training to levels comparable to those during inference.
These advancements suggest that on-device training is no longer limited to high-end devices with abundant resources. Even small embedded systems with memory capacities measured in kilobytes rather than gigabytes can participate in model training.\looseness-1
